% Exercises for Chapter 1 (sigma-algebras). Statements only: solutions are kept privately.
% Computational exercises can use the module code/src/probnotes/sigma.py (import probnotes.sigma).

\begin{exercise}[The $\sigma$-algebra generated by two sets]\label{ex:sa-two-sets}
\easy This exercise completes \noteref{exa:sa-generated}, using only the definitions of the
chapter. Let $A,B\subseteq X$, and call the four sets $A\cap B$, $A\setminus B$, $B\setminus A$ and
$(A\cup B)^c$ the \emph{cells} of $A$ and $B$.
\begin{enumerate}[label=(\alph*)]
  \item Show that the unions of cells (including the empty union) form a $\sigma$-algebra, and
    deduce that $\sigma(\{A,B\})$ consists exactly of the unions of cells.
  \item Assume that the four cells are non-empty. List the $16$ elements of $\sigma(\{A,B\})$,
    writing each of them in terms of $A$ and $B$.
  \item Assume that $A\cap B=\emptyset$ and that $A$, $B$ and $(A\cup B)^c$ are non-empty. List
    the $8$ elements of $\sigma(\{A,B\})$.
  \item Show that in general $\sigma(\{A,B\})$ has $2^k$ elements, where $k$ is the number of
    non-empty cells, and give sets $A,B\subseteq\{1,2,3,4\}$ for each $k=1,2,3,4$.
  \item More generally, show that $\sigma(\{A_1,\dots,A_m\})$ has at most $2^{2^m}$ elements, and
    that there are sets $A_1,\dots,A_m\subseteq X$ for which it has exactly $2^{2^m}$ elements if
    and only if $X$ has at least $2^m$ points.
\end{enumerate}
\end{exercise}

\begin{exercise}[Computing generated $\sigma$-algebras]\label{ex:sa-computing}
\easy This is a computational exercise; use \texttt{probnotes.sigma} (the functions
\texttt{atoms}, \texttt{generate} and \texttt{closure\_steps}) or your own code. Let
$X=\{1,\dots,6\}$.
\begin{enumerate}[label=(\alph*)]
  \item For each of the families
    \[
      \mathcal E_a=\{\{1,2\},\{2,3\}\},\quad
      \mathcal E_b=\{\{1,2,3\},\{3,4,5\},\{5,6,1\}\},\quad
      \mathcal E_c=\{\{1,2\},\{3,4\},\{5,6\}\},
    \]
    \[
      \mathcal E_d=\{\{1,2,3,4\},\{3,4,5,6\},\{1,2,5,6\}\},
    \]
    compute the atoms of $\sigma(\mathcal E)$ and its number of elements by the two methods of
    \noteref{rem:sa-atoms}, and check that they agree.
  \item Show by hand, with \noteref{prop:sa-same-generated}, that
    $\sigma(\mathcal E_c)=\sigma(\mathcal E_d)$, and confirm it numerically.
  \item Let $X$ be any finite set. A family $\mathcal E\subseteq\mathcal P(X)$ \emph{separates
    points} if for any two points $x\ne y$ of $X$ some $E\in\mathcal E$ contains exactly one of
    them. Prove that $\sigma(\mathcal E)=\mathcal P(X)$ if and only if $\mathcal E$ separates
    points.
  \item Deduce that a family of $m$ sets can generate $\mathcal P(\{1,\dots,n\})$ only if
    $n\le2^m$. For every $n\ge1$, find a family of $\lceil\log_2n\rceil$ sets that generates
    $\mathcal P(\{1,\dots,n\})$, and check your construction with code for $n=1,\dots,16$.
\end{enumerate}
\end{exercise}

\begin{exercise}[Counting $\sigma$-algebras]\label{ex:sa-counting}
\medium This is a computational exercise. Let $X_n=\{1,\dots,n\}$. An \emph{atom} of a
$\sigma$-algebra $\mathcal A$ on $X_n$ is a non-empty $A\in\mathcal A$ that contains no non-empty
element of $\mathcal A$ other than itself.
\begin{enumerate}[label=(\alph*)]
  \item Show that the atoms of $\mathcal A$ form a partition of $X_n$ and that $\mathcal A$
    consists of the unions of its atoms. Deduce that $\mathcal A\mapsto\{\text{atoms of }\mathcal
    A\}$ is a bijection between the $\sigma$-algebras on $X_n$ and the partitions of $X_n$, so that
    there are $B_n$ of them, the $n$-th Bell number. Check this by brute force over all
    $2^{2^n}$ families of subsets of $X_n$ for $n\le4$ (\texttt{count\_algebras\_bruteforce}, or
    your own code), and by listing partitions for $n\le8$.
  \item For $n\le7$ and $k=1,\dots,n$, count the $\sigma$-algebras on $X_n$ with exactly $2^k$
    elements. Show that these are the Stirling numbers of the second kind $S(n,k)$, which satisfy
    $S(n,k)=k\,S(n-1,k)+S(n-1,k-1)$, and explain this recursion in terms of $\sigma$-algebras:
    where does the point $n$ go?
  \item How many $\sigma$-algebras on $X_n$ contain the set $\{1\}$? How many contain $\{1,2\}$?
    Guess formulas in terms of Bell numbers from your counts for $n\le8$, then prove them.
\end{enumerate}
\end{exercise}

\begin{exercise}[Random generating families]\label{ex:sa-random-families}
\medium This is a computational exercise (in the style of Ross, Ch.~2). Let $X=\{1,\dots,n\}$,
and let $E_1,\dots,E_m$ be independent random subsets of $X$, each uniformly distributed on
$\mathcal P(X)$: each point of $X$ belongs to each $E_j$ with probability $\frac12$,
independently. Let $\mathcal E=\{E_1,\dots,E_m\}$.
\begin{enumerate}[label=(\alph*)]
  \item Using \exref{ex:sa-computing}(c), show that
    \[
      \Prob\big(\sigma(\mathcal E)=\mathcal P(X)\big)=\prod_{k=0}^{n-1}\Big(1-\frac{k}{2^m}\Big).
    \]
    Which classical problem of elementary probability is this?
  \item For $n=10$ and $m=4,\dots,12$, estimate this probability by simulation with $10^4$
    families for each $m$, and compare with (a). Report the estimates with $95\%$ confidence
    intervals, and your seed.
  \item For $n=10$, $100$ and $1000$, find the smallest $m$ for which the probability is at least
    $\frac12$. Using $1-x\approx e^{-x}$, explain why this $m$ is close to
    $2\log_2n-\log_2(2\ln2)$.
\end{enumerate}
\end{exercise}

\begin{exercise}[An increasing union of $\sigma$-algebras]\label{ex:sa-increasing-union}
\easy Let $X=\N$ and, for $n\in\N$, $\mathcal A_n=\sigma\big(\{0\},\{1\},\dots,\{n\}\big)$, the
$\sigma$-algebra generated by the singletons $\{0\},\dots,\{n\}$.
\begin{enumerate}[label=(\alph*)]
  \item Show that $A\in\mathcal A_n$ if and only if either $A\cap\{n+1,n+2,\dots\}=\emptyset$ or
    $\{n+1,n+2,\dots\}\subseteq A$.
  \item Show that $\mathcal A_n\subseteq\mathcal A_{n+1}$ for every $n$, and that the union of an
    increasing sequence of algebras is an algebra.
  \item Show that $\bigcup_{n\in\N}\mathcal A_n$ is the algebra of
    \noteref{exa:sa-finite-cofinite} with $X=\N$, so that it is not a $\sigma$-algebra, and that
    $\bigvee_{n\in\N}\mathcal A_n=\mathcal P(\N)$.
\end{enumerate}
Compare with \noteref{rem:sa-intersection}, where the union of two $\sigma$-algebras is not even
an algebra.
\end{exercise}

\begin{exercise}[The $\sigma$-algebra generated by the singletons]\label{ex:sa-countable-singletons}
\medium Let $X$ be an uncountable set, $\mathcal A$ the $\sigma$-algebra of countable and
cocountable sets of \noteref{exa:sa-countable-cocountable}, and $\mathcal S=\{\{x\}:x\in X\}$.
\begin{enumerate}[label=(\alph*)]
  \item Show that $\sigma(\mathcal S)=\mathcal A$.
  \item Let $X=\R$. Show that $\mathcal A\subsetneq\mathcal B(\R)$, and give a Borel set that is
    not in $\mathcal A$. So the singletons, although they are Borel sets, do not generate
    $\mathcal B(\R)$.
  \item Let $\mathcal E\subseteq\mathcal P(X)$ be any family of subsets of any set $X$. Show that
    for every $A\in\sigma(\mathcal E)$ there is an at most countable subfamily
    $\mathcal E_0\subseteq\mathcal E$ with $A\in\sigma(\mathcal E_0)$. \emph{Hint:} show that the
    sets with this property form a $\sigma$-algebra. Use this to give a second proof of
    $\sigma(\mathcal S)\subseteq\mathcal A$ in (a).
\end{enumerate}
\end{exercise}

\begin{exercise}[More generators of the Borel sets]\label{ex:sa-borel-more}
\medium You may use \noteref{prop:sa-borel-generators}, the Heine--Borel theorem, and the fact
that compact subsets of a metric space are closed. Show that
$\mathcal B(\R)$ is generated by each of the following families:
\begin{enumerate}[label=(\alph*)]
  \item the closed subsets of $\R$;
  \item the compact subsets of $\R$;
  \item the half-lines $(-\infty,q]$ with $q\in\Q$; in particular, $\mathcal B(\R)$ is generated by
    a countable family;
  \item the dyadic intervals $[k2^{-n},(k+1)2^{-n})$ with $k\in\Z$ and $n\in\N$.
\end{enumerate}
\end{exercise}

\begin{exercise}[Borel generators on a finite grid]\label{ex:sa-grid}
\medium This is a computational exercise. Let $G=\{\frac k{12}:k=0,1,\dots,12\}$ and
$D=\{0,\frac14,\frac12,\frac34,1\}$. For $i=1,\dots,8$, let $\mathcal E_i^D$ be the sets of the
family $\mathcal E_i$ of \noteref{prop:sa-borel-generators} whose endpoints lie in $D$, and let
$\mathcal G_i=\sigma\big(\{E\cap G:E\in\mathcal E_i^D\}\big)$, a $\sigma$-algebra on the finite set
$G$. The function \texttt{borel\_generator\_traces} of \texttt{probnotes.sigma} computes the sets
$E\cap G$.
\begin{enumerate}[label=(\alph*)]
  \item Compute the atoms of $\mathcal G_1,\dots,\mathcal G_8$, and their number.
  \item Show that the eight $\sigma$-algebras fall into four groups of equal ones, and that two
    $\sigma$-algebras from different groups are not comparable: neither contains the other.
    Explain the equalities $\mathcal G_5=\mathcal G_8$ and $\mathcal G_3=\mathcal G_5$ by
    identities between sets. Explain why $\mathcal G_1$ and $\mathcal G_2$ are not comparable, and
    point out the steps of the proof of \noteref{prop:sa-borel-generators} that fail on the grid.
  \item Repeat (a) with $D$ replaced by $G$, and then by
    $D'=\{\frac k{12}:k=-1,0,1,\dots,13\}$. Explain the results.
\end{enumerate}
\end{exercise}

\begin{exercise}[No $\sigma$-algebra is countably infinite]\label{ex:sa-cardinality}
\hard Let $\mathcal A$ be a $\sigma$-algebra of subsets of a non-empty set $X$. For $E\in\mathcal
A$, let $\mathcal A_E=\{A\in\mathcal A:A\subseteq E\}$.
\begin{enumerate}[label=(\alph*)]
  \item Show that if $\mathcal A$ is finite, it has $2^k$ elements for some integer $k\ge1$.
    \emph{Hint:} for $x\in X$, consider the intersection of all the sets of $\mathcal A$ that
    contain $x$.
  \item Let $E,B\in\mathcal A$ with $B\subseteq E$. Show that $A\mapsto(A\cap B,A\cap(E\setminus B))$
    is a bijection from $\mathcal A_E$ onto $\mathcal A_B\times\mathcal A_{E\setminus B}$.
  \item Suppose that $\mathcal A$ is infinite. Construct by induction sets
    $X=E_0\supseteq E_1\supseteq E_2\supseteq\cdots$ in $\mathcal A$ such that, for every $k$,
    $\mathcal A_{E_k}$ is infinite and $D_{k+1}=E_k\setminus E_{k+1}$ is non-empty.
  \item Deduce that the sets $D_k$ are pairwise disjoint, that $S\mapsto\bigcup_{k\in S}D_k$ is an
    injective map from $\mathcal P(\N^*)$ into $\mathcal A$, and hence that an infinite
    $\sigma$-algebra has at least the cardinality of the continuum. In particular, no
    $\sigma$-algebra is countably infinite.
  \item Is the same true of algebras?
\end{enumerate}
\end{exercise}
